% literature_review.tex -- Section 2 (Related work) for paper/manuscript.tex
%
% How to include it: put % literature_review.tex -- Section 2 (Related work) for paper/manuscript.tex
%
% How to include it: put % literature_review.tex -- Section 2 (Related work) for paper/manuscript.tex
%
% How to include it: put % literature_review.tex -- Section 2 (Related work) for paper/manuscript.tex
%
% How to include it: put \input{literature_review} after the Introduction and point the
% bibliography at paper/references.bib (\bibliography{references} instead of cas-refs).
% Citations use natbib author-year commands (\citep, \citet); every key is in references.bib.
%
% Drafted 2026-09-29. Each citation was checked against the source's abstract, publisher page or an
% index entry (full texts were not reachable from the drafting environment). The claim-by-claim
% audit is in paper/literature_package.md, section L. Two cited works are preprints
% (Sartori et al. 2026; Sathyamoorthy and Puri 2026) and one is an arXiv note (Opitz and Burst 2019);
% the text says so where it matters and the reference list labels them.

\section{Related work}\label{sec:related}

Four lines of earlier work shape how enzyme function predictors should be compared: the limits of homology transfer and the learned methods that followed it, the metrics and benchmarks used to rank these methods, the handling of rare and unseen EC numbers, and the reliability of prediction confidence.

\subsection{From homology transfer to learned predictors}\label{sec:related:methods}

EC numbers classify enzyme-catalysed reactions in a four-level hierarchy \citep{McDonald2023}. Because the scheme classifies reactions rather than proteins, one protein can carry several EC numbers, and unrelated proteins can catalyse the same reaction \citep{Omelchenko2010}. A long-standing way to annotate a new sequence is to transfer EC numbers from its closest annotated homologue, found with BLAST \citep{Altschul1997,Camacho2009} or with EC-specific sequence profiles such as PRIAM \citep{ClaudelRenard2003}. The reliability of this transfer falls with sequence identity. \citet{Rost2002} found that enzyme function is less conserved between homologues than had been assumed, and \citet{Tian2003} estimated that about 40\% identity is needed before more than 90\% of homologous pairs share the first three EC digits, and about 60\% before they share all four. Transfer also copies existing errors. Across 37 enzyme families, misannotation was rare in the manually reviewed Swiss-Prot but frequent in automatically annotated databases \citep{Schnoes2009}, and an experimental screen of 122 sequences from a single class, EC 1.1.3.15, implied that at least 78\% of the sequences annotated to that class were misannotated \citep{Rembeza2021}.

Machine-learning predictors were developed to reach beyond close homologues. ECPred trains an independent classifier for each EC number \citep{Dalkiran2018}, DeepEC and ProteInfer apply convolutional neural networks to the sequence \citep{Ryu2019,Sanderson2023}, and DeepFRI combines language-model features with protein structure in a graph convolutional network \citep{Gligorijevic2021}. DeepECtransformer, which uses transformer layers, proposed EC numbers for 464 unannotated \textit{Escherichia coli} genes and confirmed three of them experimentally \citep{Kim2023}, while SOLVE showed that interpretable ensembles built on short sequence fragments can compete with deep models \citep{Dhibar2025}. CLEAN \citep{Yu2023} takes embeddings from the protein language model ESM-1b \citep{Rives2021}, trains a contrastive embedding on them and assigns the EC numbers whose cluster centres lie nearest to the query. It was reported to be more accurate and sensitive than BLASTp and has since been used as a reference method in independent benchmarks \citep{Davoudi2026,Sartori2026}. Later methods introduce predicted contact maps \citep[CLEAN-Contact;][]{YangY2024}, residue-level detection \citep[ProtDETR;][]{YangZ2025}, hierarchy-aware contrastive losses with structural information \citep[EnzymeHunter;][]{Cao2026}, the prediction of EC numbers absent from the training data \citep[EnzPlacer;][]{Ma2026} and a hierarchical transformer \citep[HIT-EC;][]{Dumontet2026}.

Independent evaluations have been less favourable than the original reports. An expert analysis of DeepECtransformer predictions for more than 450 \textit{E.~coli} enzymes of unknown function concluded that current models mostly fail to make novel predictions and make basic logic errors that human annotators avoid \citep{deCrecyLagard2025}. In a 2026 preprint, CLEAN and the homology tool DIAMOND still assigned the original EC number to more than 90\% of weakly perturbed decoys whose active sites had been deliberately disrupted \citep{Sartori2026}. Whether learned embeddings outperform homology search at all is also contested. A comparison of ESM2, ESM-1b and ProtBERT classifiers with BLASTp found BLASTp marginally better overall, although the two approaches were complementary: the language models, ESM2 in particular, predicted some EC numbers better and were more useful for enzymes with less than 25\% identity to the reference database \citep{Capela2025}. Another 2026 preprint, which matched the BLASTp database to each model's training data, found language-model classifiers level with BLASTp on in-distribution proteins and well ahead of it on evolutionarily distant eukaryotic enzymes \citep{Sathyamoorthy2026}. The outcome therefore appears to depend on the identity profile of the test set and on how the homology baseline is configured. For this reason we compare CLEAN with a standard BLASTp search against CLEAN's own training sequences and report the results by identity stratum.

\subsection{Metrics and benchmarks}\label{sec:related:metrics}

Community assessments of protein function prediction evaluate methods prospectively. In the CAFA challenges, methods predict functions for proteins that lack experimental annotations and are scored on the annotations that accumulate afterwards \citep{Radivojac2013,Jiang2016,Zhou2019}; LAFA now runs this time-delayed evaluation continuously \citep{Phan2026}. Two lessons from the evaluation of Gene Ontology predictors apply directly to EC prediction. First, the metric can decide the ranking: \citet{Plyusnin2019} compared 37 evaluation metrics for Gene Ontology classifiers and found that some of them gave erroneously high scores to deliberately constructed false-positive annotation sets. Second, the ground truth is incomplete, so a prediction scored as a false positive may be a true function that has not yet been annotated \citep{Jiang2014}.

EC predictors, by contrast, are usually compared with standard multi-label precision, recall and F1. These can be averaged over proteins, averaged over labels, or pooled over all decisions by micro-averaging \citep{Zhang2014}. These choices are not interchangeable. Measures differ in which changes to the confusion matrix they are invariant to \citep{Sokolova2009}, and even the two formulas in use for macro-F1 can rank the same classifiers differently \citep{Opitz2019}. EC studies commonly average over EC classes with weights proportional to class support; EC-Bench, for example, adopts the support-weighted F1 to account for class imbalance \citep{Davoudi2026}. Recent benchmarks have also standardised tasks and controlled leakage. CARE tests EC classification on sequences with less than 30\% or 30--50\% identity to the training data and adds the retrieval of EC numbers from reactions \citep{YangJ2024}, and EC-Bench compares ten existing methods on exact EC prediction, EC completion and EC recommendation under uniform assumptions \citep{Davoudi2026}. Temporal hold-outs have a longer history in EC prediction. ECPred was tested on a small set of proteins that received an EC number only in a later Swiss-Prot release and reached an F1 of 0.14 at the substrate level; its authors noted that scores on such proteins can differ greatly from published values because these cases are very difficult \citep{Dalkiran2018}. These efforts are in line with general guidance on validation and data leakage in biological machine learning \citep{Walsh2021,Kapoor2023}. To our knowledge, however, no study has asked whether the averaging used to score EC predictors counts every false positive. We show that the support-weighted form does not.

\subsection{Rare and unseen EC numbers}\label{sec:related:rare}

EC labels are imbalanced \citep{Davoudi2026}, and several recent methods target the rare classes. HIT-EC was motivated by the low performance of existing models on under-represented EC numbers and by incomplete annotations \citep{Dumontet2026}, and EnzPlacer predicts EC numbers that are absent from the training set altogether \citep{Ma2026}. A closed-set predictor such as CLEAN can only output EC numbers seen in training. The test labels it has never seen therefore cap its attainable recall, and a single headline score on a time-split test set mixes proteins with unseen EC numbers and proteins whose EC numbers were seen. For the seen classes, the inference rule matters. CLEAN scores each EC number by the distance between the query and the mean embedding of that EC number's training sequences. This is the prototype classifier of few-shot learning \citep{Snell2017}, applied to an embedding trained with a supervised contrastive objective \citep{Khosla2020}. In long-tailed recognition, prototype classifiers such as the nearest class mean have been proposed to counter the bias of standard classifiers towards frequent classes \citep{Sharma2023}. To our knowledge, centroid and nearest-neighbour rules have not been compared within one frozen EC embedding space, so it is not known which of them serves rare EC numbers better.

\subsection{Confidence, calibration and guarantees}\label{sec:related:calibration}

For a curator, a predicted EC number is most useful when its confidence can be trusted. Modern neural networks are often poorly calibrated \citep{Guo2017}, and their scores can be mapped to calibrated probabilities after training, for example by isotonic regression \citep{Zadrozny2002}. Selective classification goes one step further: the model abstains when its confidence is low, and the error rate of the predictions it keeps is guaranteed with high probability \citep{Geifman2017}. For EC prediction, CLEAN expresses confidence by fitting a two-component Gaussian mixture to embedding distances \citep{Yu2023}, and ECPICK uses evidential deep learning \citep{Han2024}. Conformal methods add distribution-free guarantees. CPEC limits the false discovery rate of EC annotations \citep{DingK2024}, and \citet{Boger2025} applied conformal methods to protein search with user-chosen risk guarantees, including search with CLEAN embeddings. These guarantees hold on average over test proteins. A guarantee for each EC class, or for each tier of rare classes, needs calibration data from that class, and class-conditional conformal prediction becomes difficult when there are many classes with few examples each \citep{DingT2023}. Exact binomial confidence limits \citep{Clopper1934} determine the minimum number of calibration proteins that such a guarantee needs. To our knowledge, nobody has measured how far the rare EC tiers are from that minimum, or how long curation would take to close the gap.

\subsection{Open questions}\label{sec:related:open}

The field therefore has strong predictors and increasingly careful benchmarks, but three questions remain open for EC prediction. Do the standard scores count all false positives? How much of a temporal test set lies beyond the reach of a closed-set model, and how does a frozen contrastive embedding compare with homology search on such a set? And how much calibration data do rare EC numbers need before their predictions can be certified? This study addresses these questions in turn.
 after the Introduction and point the
% bibliography at paper/references.bib (\bibliography{references} instead of cas-refs).
% Citations use natbib author-year commands (\citep, \citet); every key is in references.bib.
%
% Drafted 2026-09-29. Each citation was checked against the source's abstract, publisher page or an
% index entry (full texts were not reachable from the drafting environment). The claim-by-claim
% audit is in paper/literature_package.md, section L. Two cited works are preprints
% (Sartori et al. 2026; Sathyamoorthy and Puri 2026) and one is an arXiv note (Opitz and Burst 2019);
% the text says so where it matters and the reference list labels them.

\section{Related work}\label{sec:related}

Four lines of earlier work shape how enzyme function predictors should be compared: the limits of homology transfer and the learned methods that followed it, the metrics and benchmarks used to rank these methods, the handling of rare and unseen EC numbers, and the reliability of prediction confidence.

\subsection{From homology transfer to learned predictors}\label{sec:related:methods}

EC numbers classify enzyme-catalysed reactions in a four-level hierarchy \citep{McDonald2023}. Because the scheme classifies reactions rather than proteins, one protein can carry several EC numbers, and unrelated proteins can catalyse the same reaction \citep{Omelchenko2010}. A long-standing way to annotate a new sequence is to transfer EC numbers from its closest annotated homologue, found with BLAST \citep{Altschul1997,Camacho2009} or with EC-specific sequence profiles such as PRIAM \citep{ClaudelRenard2003}. The reliability of this transfer falls with sequence identity. \citet{Rost2002} found that enzyme function is less conserved between homologues than had been assumed, and \citet{Tian2003} estimated that about 40\% identity is needed before more than 90\% of homologous pairs share the first three EC digits, and about 60\% before they share all four. Transfer also copies existing errors. Across 37 enzyme families, misannotation was rare in the manually reviewed Swiss-Prot but frequent in automatically annotated databases \citep{Schnoes2009}, and an experimental screen of 122 sequences from a single class, EC 1.1.3.15, implied that at least 78\% of the sequences annotated to that class were misannotated \citep{Rembeza2021}.

Machine-learning predictors were developed to reach beyond close homologues. ECPred trains an independent classifier for each EC number \citep{Dalkiran2018}, DeepEC and ProteInfer apply convolutional neural networks to the sequence \citep{Ryu2019,Sanderson2023}, and DeepFRI combines language-model features with protein structure in a graph convolutional network \citep{Gligorijevic2021}. DeepECtransformer, which uses transformer layers, proposed EC numbers for 464 unannotated \textit{Escherichia coli} genes and confirmed three of them experimentally \citep{Kim2023}, while SOLVE showed that interpretable ensembles built on short sequence fragments can compete with deep models \citep{Dhibar2025}. CLEAN \citep{Yu2023} takes embeddings from the protein language model ESM-1b \citep{Rives2021}, trains a contrastive embedding on them and assigns the EC numbers whose cluster centres lie nearest to the query. It was reported to be more accurate and sensitive than BLASTp and has since been used as a reference method in independent benchmarks \citep{Davoudi2026,Sartori2026}. Later methods introduce predicted contact maps \citep[CLEAN-Contact;][]{YangY2024}, residue-level detection \citep[ProtDETR;][]{YangZ2025}, hierarchy-aware contrastive losses with structural information \citep[EnzymeHunter;][]{Cao2026}, the prediction of EC numbers absent from the training data \citep[EnzPlacer;][]{Ma2026} and a hierarchical transformer \citep[HIT-EC;][]{Dumontet2026}.

Independent evaluations have been less favourable than the original reports. An expert analysis of DeepECtransformer predictions for more than 450 \textit{E.~coli} enzymes of unknown function concluded that current models mostly fail to make novel predictions and make basic logic errors that human annotators avoid \citep{deCrecyLagard2025}. In a 2026 preprint, CLEAN and the homology tool DIAMOND still assigned the original EC number to more than 90\% of weakly perturbed decoys whose active sites had been deliberately disrupted \citep{Sartori2026}. Whether learned embeddings outperform homology search at all is also contested. A comparison of ESM2, ESM-1b and ProtBERT classifiers with BLASTp found BLASTp marginally better overall, although the two approaches were complementary: the language models, ESM2 in particular, predicted some EC numbers better and were more useful for enzymes with less than 25\% identity to the reference database \citep{Capela2025}. Another 2026 preprint, which matched the BLASTp database to each model's training data, found language-model classifiers level with BLASTp on in-distribution proteins and well ahead of it on evolutionarily distant eukaryotic enzymes \citep{Sathyamoorthy2026}. The outcome therefore appears to depend on the identity profile of the test set and on how the homology baseline is configured. For this reason we compare CLEAN with a standard BLASTp search against CLEAN's own training sequences and report the results by identity stratum.

\subsection{Metrics and benchmarks}\label{sec:related:metrics}

Community assessments of protein function prediction evaluate methods prospectively. In the CAFA challenges, methods predict functions for proteins that lack experimental annotations and are scored on the annotations that accumulate afterwards \citep{Radivojac2013,Jiang2016,Zhou2019}; LAFA now runs this time-delayed evaluation continuously \citep{Phan2026}. Two lessons from the evaluation of Gene Ontology predictors apply directly to EC prediction. First, the metric can decide the ranking: \citet{Plyusnin2019} compared 37 evaluation metrics for Gene Ontology classifiers and found that some of them gave erroneously high scores to deliberately constructed false-positive annotation sets. Second, the ground truth is incomplete, so a prediction scored as a false positive may be a true function that has not yet been annotated \citep{Jiang2014}.

EC predictors, by contrast, are usually compared with standard multi-label precision, recall and F1. These can be averaged over proteins, averaged over labels, or pooled over all decisions by micro-averaging \citep{Zhang2014}. These choices are not interchangeable. Measures differ in which changes to the confusion matrix they are invariant to \citep{Sokolova2009}, and even the two formulas in use for macro-F1 can rank the same classifiers differently \citep{Opitz2019}. EC studies commonly average over EC classes with weights proportional to class support; EC-Bench, for example, adopts the support-weighted F1 to account for class imbalance \citep{Davoudi2026}. Recent benchmarks have also standardised tasks and controlled leakage. CARE tests EC classification on sequences with less than 30\% or 30--50\% identity to the training data and adds the retrieval of EC numbers from reactions \citep{YangJ2024}, and EC-Bench compares ten existing methods on exact EC prediction, EC completion and EC recommendation under uniform assumptions \citep{Davoudi2026}. Temporal hold-outs have a longer history in EC prediction. ECPred was tested on a small set of proteins that received an EC number only in a later Swiss-Prot release and reached an F1 of 0.14 at the substrate level; its authors noted that scores on such proteins can differ greatly from published values because these cases are very difficult \citep{Dalkiran2018}. These efforts are in line with general guidance on validation and data leakage in biological machine learning \citep{Walsh2021,Kapoor2023}. To our knowledge, however, no study has asked whether the averaging used to score EC predictors counts every false positive. We show that the support-weighted form does not.

\subsection{Rare and unseen EC numbers}\label{sec:related:rare}

EC labels are imbalanced \citep{Davoudi2026}, and several recent methods target the rare classes. HIT-EC was motivated by the low performance of existing models on under-represented EC numbers and by incomplete annotations \citep{Dumontet2026}, and EnzPlacer predicts EC numbers that are absent from the training set altogether \citep{Ma2026}. A closed-set predictor such as CLEAN can only output EC numbers seen in training. The test labels it has never seen therefore cap its attainable recall, and a single headline score on a time-split test set mixes proteins with unseen EC numbers and proteins whose EC numbers were seen. For the seen classes, the inference rule matters. CLEAN scores each EC number by the distance between the query and the mean embedding of that EC number's training sequences. This is the prototype classifier of few-shot learning \citep{Snell2017}, applied to an embedding trained with a supervised contrastive objective \citep{Khosla2020}. In long-tailed recognition, prototype classifiers such as the nearest class mean have been proposed to counter the bias of standard classifiers towards frequent classes \citep{Sharma2023}. To our knowledge, centroid and nearest-neighbour rules have not been compared within one frozen EC embedding space, so it is not known which of them serves rare EC numbers better.

\subsection{Confidence, calibration and guarantees}\label{sec:related:calibration}

For a curator, a predicted EC number is most useful when its confidence can be trusted. Modern neural networks are often poorly calibrated \citep{Guo2017}, and their scores can be mapped to calibrated probabilities after training, for example by isotonic regression \citep{Zadrozny2002}. Selective classification goes one step further: the model abstains when its confidence is low, and the error rate of the predictions it keeps is guaranteed with high probability \citep{Geifman2017}. For EC prediction, CLEAN expresses confidence by fitting a two-component Gaussian mixture to embedding distances \citep{Yu2023}, and ECPICK uses evidential deep learning \citep{Han2024}. Conformal methods add distribution-free guarantees. CPEC limits the false discovery rate of EC annotations \citep{DingK2024}, and \citet{Boger2025} applied conformal methods to protein search with user-chosen risk guarantees, including search with CLEAN embeddings. These guarantees hold on average over test proteins. A guarantee for each EC class, or for each tier of rare classes, needs calibration data from that class, and class-conditional conformal prediction becomes difficult when there are many classes with few examples each \citep{DingT2023}. Exact binomial confidence limits \citep{Clopper1934} determine the minimum number of calibration proteins that such a guarantee needs. To our knowledge, nobody has measured how far the rare EC tiers are from that minimum, or how long curation would take to close the gap.

\subsection{Open questions}\label{sec:related:open}

The field therefore has strong predictors and increasingly careful benchmarks, but three questions remain open for EC prediction. Do the standard scores count all false positives? How much of a temporal test set lies beyond the reach of a closed-set model, and how does a frozen contrastive embedding compare with homology search on such a set? And how much calibration data do rare EC numbers need before their predictions can be certified? This study addresses these questions in turn.
 after the Introduction and point the
% bibliography at paper/references.bib (\bibliography{references} instead of cas-refs).
% Citations use natbib author-year commands (\citep, \citet); every key is in references.bib.
%
% Drafted 2026-09-29. Each citation was checked against the source's abstract, publisher page or an
% index entry (full texts were not reachable from the drafting environment). The claim-by-claim
% audit is in paper/literature_package.md, section L. Two cited works are preprints
% (Sartori et al. 2026; Sathyamoorthy and Puri 2026) and one is an arXiv note (Opitz and Burst 2019);
% the text says so where it matters and the reference list labels them.

\section{Related work}\label{sec:related}

Four lines of earlier work shape how enzyme function predictors should be compared: the limits of homology transfer and the learned methods that followed it, the metrics and benchmarks used to rank these methods, the handling of rare and unseen EC numbers, and the reliability of prediction confidence.

\subsection{From homology transfer to learned predictors}\label{sec:related:methods}

EC numbers classify enzyme-catalysed reactions in a four-level hierarchy \citep{McDonald2023}. Because the scheme classifies reactions rather than proteins, one protein can carry several EC numbers, and unrelated proteins can catalyse the same reaction \citep{Omelchenko2010}. A long-standing way to annotate a new sequence is to transfer EC numbers from its closest annotated homologue, found with BLAST \citep{Altschul1997,Camacho2009} or with EC-specific sequence profiles such as PRIAM \citep{ClaudelRenard2003}. The reliability of this transfer falls with sequence identity. \citet{Rost2002} found that enzyme function is less conserved between homologues than had been assumed, and \citet{Tian2003} estimated that about 40\% identity is needed before more than 90\% of homologous pairs share the first three EC digits, and about 60\% before they share all four. Transfer also copies existing errors. Across 37 enzyme families, misannotation was rare in the manually reviewed Swiss-Prot but frequent in automatically annotated databases \citep{Schnoes2009}, and an experimental screen of 122 sequences from a single class, EC 1.1.3.15, implied that at least 78\% of the sequences annotated to that class were misannotated \citep{Rembeza2021}.

Machine-learning predictors were developed to reach beyond close homologues. ECPred trains an independent classifier for each EC number \citep{Dalkiran2018}, DeepEC and ProteInfer apply convolutional neural networks to the sequence \citep{Ryu2019,Sanderson2023}, and DeepFRI combines language-model features with protein structure in a graph convolutional network \citep{Gligorijevic2021}. DeepECtransformer, which uses transformer layers, proposed EC numbers for 464 unannotated \textit{Escherichia coli} genes and confirmed three of them experimentally \citep{Kim2023}, while SOLVE showed that interpretable ensembles built on short sequence fragments can compete with deep models \citep{Dhibar2025}. CLEAN \citep{Yu2023} takes embeddings from the protein language model ESM-1b \citep{Rives2021}, trains a contrastive embedding on them and assigns the EC numbers whose cluster centres lie nearest to the query. It was reported to be more accurate and sensitive than BLASTp and has since been used as a reference method in independent benchmarks \citep{Davoudi2026,Sartori2026}. Later methods introduce predicted contact maps \citep[CLEAN-Contact;][]{YangY2024}, residue-level detection \citep[ProtDETR;][]{YangZ2025}, hierarchy-aware contrastive losses with structural information \citep[EnzymeHunter;][]{Cao2026}, the prediction of EC numbers absent from the training data \citep[EnzPlacer;][]{Ma2026} and a hierarchical transformer \citep[HIT-EC;][]{Dumontet2026}.

Independent evaluations have been less favourable than the original reports. An expert analysis of DeepECtransformer predictions for more than 450 \textit{E.~coli} enzymes of unknown function concluded that current models mostly fail to make novel predictions and make basic logic errors that human annotators avoid \citep{deCrecyLagard2025}. In a 2026 preprint, CLEAN and the homology tool DIAMOND still assigned the original EC number to more than 90\% of weakly perturbed decoys whose active sites had been deliberately disrupted \citep{Sartori2026}. Whether learned embeddings outperform homology search at all is also contested. A comparison of ESM2, ESM-1b and ProtBERT classifiers with BLASTp found BLASTp marginally better overall, although the two approaches were complementary: the language models, ESM2 in particular, predicted some EC numbers better and were more useful for enzymes with less than 25\% identity to the reference database \citep{Capela2025}. Another 2026 preprint, which matched the BLASTp database to each model's training data, found language-model classifiers level with BLASTp on in-distribution proteins and well ahead of it on evolutionarily distant eukaryotic enzymes \citep{Sathyamoorthy2026}. The outcome therefore appears to depend on the identity profile of the test set and on how the homology baseline is configured. For this reason we compare CLEAN with a standard BLASTp search against CLEAN's own training sequences and report the results by identity stratum.

\subsection{Metrics and benchmarks}\label{sec:related:metrics}

Community assessments of protein function prediction evaluate methods prospectively. In the CAFA challenges, methods predict functions for proteins that lack experimental annotations and are scored on the annotations that accumulate afterwards \citep{Radivojac2013,Jiang2016,Zhou2019}; LAFA now runs this time-delayed evaluation continuously \citep{Phan2026}. Two lessons from the evaluation of Gene Ontology predictors apply directly to EC prediction. First, the metric can decide the ranking: \citet{Plyusnin2019} compared 37 evaluation metrics for Gene Ontology classifiers and found that some of them gave erroneously high scores to deliberately constructed false-positive annotation sets. Second, the ground truth is incomplete, so a prediction scored as a false positive may be a true function that has not yet been annotated \citep{Jiang2014}.

EC predictors, by contrast, are usually compared with standard multi-label precision, recall and F1. These can be averaged over proteins, averaged over labels, or pooled over all decisions by micro-averaging \citep{Zhang2014}. These choices are not interchangeable. Measures differ in which changes to the confusion matrix they are invariant to \citep{Sokolova2009}, and even the two formulas in use for macro-F1 can rank the same classifiers differently \citep{Opitz2019}. EC studies commonly average over EC classes with weights proportional to class support; EC-Bench, for example, adopts the support-weighted F1 to account for class imbalance \citep{Davoudi2026}. Recent benchmarks have also standardised tasks and controlled leakage. CARE tests EC classification on sequences with less than 30\% or 30--50\% identity to the training data and adds the retrieval of EC numbers from reactions \citep{YangJ2024}, and EC-Bench compares ten existing methods on exact EC prediction, EC completion and EC recommendation under uniform assumptions \citep{Davoudi2026}. Temporal hold-outs have a longer history in EC prediction. ECPred was tested on a small set of proteins that received an EC number only in a later Swiss-Prot release and reached an F1 of 0.14 at the substrate level; its authors noted that scores on such proteins can differ greatly from published values because these cases are very difficult \citep{Dalkiran2018}. These efforts are in line with general guidance on validation and data leakage in biological machine learning \citep{Walsh2021,Kapoor2023}. To our knowledge, however, no study has asked whether the averaging used to score EC predictors counts every false positive. We show that the support-weighted form does not.

\subsection{Rare and unseen EC numbers}\label{sec:related:rare}

EC labels are imbalanced \citep{Davoudi2026}, and several recent methods target the rare classes. HIT-EC was motivated by the low performance of existing models on under-represented EC numbers and by incomplete annotations \citep{Dumontet2026}, and EnzPlacer predicts EC numbers that are absent from the training set altogether \citep{Ma2026}. A closed-set predictor such as CLEAN can only output EC numbers seen in training. The test labels it has never seen therefore cap its attainable recall, and a single headline score on a time-split test set mixes proteins with unseen EC numbers and proteins whose EC numbers were seen. For the seen classes, the inference rule matters. CLEAN scores each EC number by the distance between the query and the mean embedding of that EC number's training sequences. This is the prototype classifier of few-shot learning \citep{Snell2017}, applied to an embedding trained with a supervised contrastive objective \citep{Khosla2020}. In long-tailed recognition, prototype classifiers such as the nearest class mean have been proposed to counter the bias of standard classifiers towards frequent classes \citep{Sharma2023}. To our knowledge, centroid and nearest-neighbour rules have not been compared within one frozen EC embedding space, so it is not known which of them serves rare EC numbers better.

\subsection{Confidence, calibration and guarantees}\label{sec:related:calibration}

For a curator, a predicted EC number is most useful when its confidence can be trusted. Modern neural networks are often poorly calibrated \citep{Guo2017}, and their scores can be mapped to calibrated probabilities after training, for example by isotonic regression \citep{Zadrozny2002}. Selective classification goes one step further: the model abstains when its confidence is low, and the error rate of the predictions it keeps is guaranteed with high probability \citep{Geifman2017}. For EC prediction, CLEAN expresses confidence by fitting a two-component Gaussian mixture to embedding distances \citep{Yu2023}, and ECPICK uses evidential deep learning \citep{Han2024}. Conformal methods add distribution-free guarantees. CPEC limits the false discovery rate of EC annotations \citep{DingK2024}, and \citet{Boger2025} applied conformal methods to protein search with user-chosen risk guarantees, including search with CLEAN embeddings. These guarantees hold on average over test proteins. A guarantee for each EC class, or for each tier of rare classes, needs calibration data from that class, and class-conditional conformal prediction becomes difficult when there are many classes with few examples each \citep{DingT2023}. Exact binomial confidence limits \citep{Clopper1934} determine the minimum number of calibration proteins that such a guarantee needs. To our knowledge, nobody has measured how far the rare EC tiers are from that minimum, or how long curation would take to close the gap.

\subsection{Open questions}\label{sec:related:open}

The field therefore has strong predictors and increasingly careful benchmarks, but three questions remain open for EC prediction. Do the standard scores count all false positives? How much of a temporal test set lies beyond the reach of a closed-set model, and how does a frozen contrastive embedding compare with homology search on such a set? And how much calibration data do rare EC numbers need before their predictions can be certified? This study addresses these questions in turn.
 after the Introduction and point the
% bibliography at paper/references.bib (\bibliography{references} instead of cas-refs).
% Citations use natbib author-year commands (\citep, \citet); every key is in references.bib.
%
% Drafted 2026-09-29. Each citation was checked against the source's abstract, publisher page or an
% index entry (full texts were not reachable from the drafting environment). The claim-by-claim
% audit is in paper/literature_package.md, section L. Two cited works are preprints
% (Sartori et al. 2026; Sathyamoorthy and Puri 2026) and one is an arXiv note (Opitz and Burst 2019);
% the text says so where it matters and the reference list labels them.

\section{Related work}\label{sec:related}

Four lines of earlier work shape how enzyme function predictors should be compared: the limits of homology transfer and the learned methods that followed it, the metrics and benchmarks used to rank these methods, the handling of rare and unseen EC numbers, and the reliability of prediction confidence.

\subsection{From homology transfer to learned predictors}\label{sec:related:methods}

EC numbers classify enzyme-catalysed reactions in a four-level hierarchy \citep{McDonald2023}. Because the scheme classifies reactions rather than proteins, one protein can carry several EC numbers, and unrelated proteins can catalyse the same reaction \citep{Omelchenko2010}. A long-standing way to annotate a new sequence is to transfer EC numbers from its closest annotated homologue, found with BLAST \citep{Altschul1997,Camacho2009} or with EC-specific sequence profiles such as PRIAM \citep{ClaudelRenard2003}. The reliability of this transfer falls with sequence identity. \citet{Rost2002} found that enzyme function is less conserved between homologues than had been assumed, and \citet{Tian2003} estimated that about 40\% identity is needed before more than 90\% of homologous pairs share the first three EC digits, and about 60\% before they share all four. Transfer also copies existing errors. Across 37 enzyme families, misannotation was rare in the manually reviewed Swiss-Prot but frequent in automatically annotated databases \citep{Schnoes2009}, and an experimental screen of 122 sequences from a single class, EC 1.1.3.15, implied that at least 78\% of the sequences annotated to that class were misannotated \citep{Rembeza2021}.

Machine-learning predictors were developed to reach beyond close homologues. ECPred trains an independent classifier for each EC number \citep{Dalkiran2018}, DeepEC and ProteInfer apply convolutional neural networks to the sequence \citep{Ryu2019,Sanderson2023}, and DeepFRI combines language-model features with protein structure in a graph convolutional network \citep{Gligorijevic2021}. DeepECtransformer, which uses transformer layers, proposed EC numbers for 464 unannotated \textit{Escherichia coli} genes and confirmed three of them experimentally \citep{Kim2023}, while SOLVE showed that interpretable ensembles built on short sequence fragments can compete with deep models \citep{Dhibar2025}. CLEAN \citep{Yu2023} takes embeddings from the protein language model ESM-1b \citep{Rives2021}, trains a contrastive embedding on them and assigns the EC numbers whose cluster centres lie nearest to the query. It was reported to be more accurate and sensitive than BLASTp and has since been used as a reference method in independent benchmarks \citep{Davoudi2026,Sartori2026}. Later methods introduce predicted contact maps \citep[CLEAN-Contact;][]{YangY2024}, residue-level detection \citep[ProtDETR;][]{YangZ2025}, hierarchy-aware contrastive losses with structural information \citep[EnzymeHunter;][]{Cao2026}, the prediction of EC numbers absent from the training data \citep[EnzPlacer;][]{Ma2026} and a hierarchical transformer \citep[HIT-EC;][]{Dumontet2026}.

Independent evaluations have been less favourable than the original reports. An expert analysis of DeepECtransformer predictions for more than 450 \textit{E.~coli} enzymes of unknown function concluded that current models mostly fail to make novel predictions and make basic logic errors that human annotators avoid \citep{deCrecyLagard2025}. In a 2026 preprint, CLEAN and the homology tool DIAMOND still assigned the original EC number to more than 90\% of weakly perturbed decoys whose active sites had been deliberately disrupted \citep{Sartori2026}. Whether learned embeddings outperform homology search at all is also contested. A comparison of ESM2, ESM-1b and ProtBERT classifiers with BLASTp found BLASTp marginally better overall, although the two approaches were complementary: the language models, ESM2 in particular, predicted some EC numbers better and were more useful for enzymes with less than 25\% identity to the reference database \citep{Capela2025}. Another 2026 preprint, which matched the BLASTp database to each model's training data, found language-model classifiers level with BLASTp on in-distribution proteins and well ahead of it on evolutionarily distant eukaryotic enzymes \citep{Sathyamoorthy2026}. The outcome therefore appears to depend on the identity profile of the test set and on how the homology baseline is configured. For this reason we compare CLEAN with a standard BLASTp search against CLEAN's own training sequences and report the results by identity stratum.

\subsection{Metrics and benchmarks}\label{sec:related:metrics}

Community assessments of protein function prediction evaluate methods prospectively. In the CAFA challenges, methods predict functions for proteins that lack experimental annotations and are scored on the annotations that accumulate afterwards \citep{Radivojac2013,Jiang2016,Zhou2019}; LAFA now runs this time-delayed evaluation continuously \citep{Phan2026}. Two lessons from the evaluation of Gene Ontology predictors apply directly to EC prediction. First, the metric can decide the ranking: \citet{Plyusnin2019} compared 37 evaluation metrics for Gene Ontology classifiers and found that some of them gave erroneously high scores to deliberately constructed false-positive annotation sets. Second, the ground truth is incomplete, so a prediction scored as a false positive may be a true function that has not yet been annotated \citep{Jiang2014}.

EC predictors, by contrast, are usually compared with standard multi-label precision, recall and F1. These can be averaged over proteins, averaged over labels, or pooled over all decisions by micro-averaging \citep{Zhang2014}. These choices are not interchangeable. Measures differ in which changes to the confusion matrix they are invariant to \citep{Sokolova2009}, and even the two formulas in use for macro-F1 can rank the same classifiers differently \citep{Opitz2019}. EC studies commonly average over EC classes with weights proportional to class support; EC-Bench, for example, adopts the support-weighted F1 to account for class imbalance \citep{Davoudi2026}. Recent benchmarks have also standardised tasks and controlled leakage. CARE tests EC classification on sequences with less than 30\% or 30--50\% identity to the training data and adds the retrieval of EC numbers from reactions \citep{YangJ2024}, and EC-Bench compares ten existing methods on exact EC prediction, EC completion and EC recommendation under uniform assumptions \citep{Davoudi2026}. Temporal hold-outs have a longer history in EC prediction. ECPred was tested on a small set of proteins that received an EC number only in a later Swiss-Prot release and reached an F1 of 0.14 at the substrate level; its authors noted that scores on such proteins can differ greatly from published values because these cases are very difficult \citep{Dalkiran2018}. These efforts are in line with general guidance on validation and data leakage in biological machine learning \citep{Walsh2021,Kapoor2023}. To our knowledge, however, no study has asked whether the averaging used to score EC predictors counts every false positive. We show that the support-weighted form does not.

\subsection{Rare and unseen EC numbers}\label{sec:related:rare}

EC labels are imbalanced \citep{Davoudi2026}, and several recent methods target the rare classes. HIT-EC was motivated by the low performance of existing models on under-represented EC numbers and by incomplete annotations \citep{Dumontet2026}, and EnzPlacer predicts EC numbers that are absent from the training set altogether \citep{Ma2026}. A closed-set predictor such as CLEAN can only output EC numbers seen in training. The test labels it has never seen therefore cap its attainable recall, and a single headline score on a time-split test set mixes proteins with unseen EC numbers and proteins whose EC numbers were seen. For the seen classes, the inference rule matters. CLEAN scores each EC number by the distance between the query and the mean embedding of that EC number's training sequences. This is the prototype classifier of few-shot learning \citep{Snell2017}, applied to an embedding trained with a supervised contrastive objective \citep{Khosla2020}. In long-tailed recognition, prototype classifiers such as the nearest class mean have been proposed to counter the bias of standard classifiers towards frequent classes \citep{Sharma2023}. To our knowledge, centroid and nearest-neighbour rules have not been compared within one frozen EC embedding space, so it is not known which of them serves rare EC numbers better.

\subsection{Confidence, calibration and guarantees}\label{sec:related:calibration}

For a curator, a predicted EC number is most useful when its confidence can be trusted. Modern neural networks are often poorly calibrated \citep{Guo2017}, and their scores can be mapped to calibrated probabilities after training, for example by isotonic regression \citep{Zadrozny2002}. Selective classification goes one step further: the model abstains when its confidence is low, and the error rate of the predictions it keeps is guaranteed with high probability \citep{Geifman2017}. For EC prediction, CLEAN expresses confidence by fitting a two-component Gaussian mixture to embedding distances \citep{Yu2023}, and ECPICK uses evidential deep learning \citep{Han2024}. Conformal methods add distribution-free guarantees. CPEC limits the false discovery rate of EC annotations \citep{DingK2024}, and \citet{Boger2025} applied conformal methods to protein search with user-chosen risk guarantees, including search with CLEAN embeddings. These guarantees hold on average over test proteins. A guarantee for each EC class, or for each tier of rare classes, needs calibration data from that class, and class-conditional conformal prediction becomes difficult when there are many classes with few examples each \citep{DingT2023}. Exact binomial confidence limits \citep{Clopper1934} determine the minimum number of calibration proteins that such a guarantee needs. To our knowledge, nobody has measured how far the rare EC tiers are from that minimum, or how long curation would take to close the gap.

\subsection{Open questions}\label{sec:related:open}

The field therefore has strong predictors and increasingly careful benchmarks, but three questions remain open for EC prediction. Do the standard scores count all false positives? How much of a temporal test set lies beyond the reach of a closed-set model, and how does a frozen contrastive embedding compare with homology search on such a set? And how much calibration data do rare EC numbers need before their predictions can be certified? This study addresses these questions in turn.
